% !TeX root = ../papers/nc-p3.tex
\section{曲線グラフでの長方形切断と四つの純関係の座標枠}
\label{sec:curve-graph-pure-frames}
\begin{theorem}\label{thm:curve-graph-pure-frames}
四辺の全射影ファイバーが空または一点の標準有限Frobenius窓で、共通グラフ $X$ が一次元とする。
一辺の点座標でベクトル次元 $a,b$ の任意の二部分空間へ制限した関係の階数 $r$ は
\[
        r\ge a+b-3+\mathbf1_{a\ne b}\qquad(1\le a,b\le4)
\]
を満たす。各辺の相異なる純関係は高々四点で、どの四点以下についても両側の文字は
それぞれ線型独立である。元代数の全次数の正則性は仮定しない。
\end{theorem}
\begin{proof}[制限した辺の両次数]
$\mathbb P^{a-1}\times\mathbb P^{b-1}$ 上で $r$ 個の双一次式を取る。
交叉は共通グラフの部分スキームなので次元1以下であり、Segre超平面の次元定理から
$r\ge a+b-3$ である。等号で積の次元が1以上なら、純一次元の完全交叉となる。
両射影への次数は $(H_1+H_2)^{a+b-3}$ の交叉数から
$\binom{a+b-3}{a-2},\binom{a+b-3}{a-1}$ である（範囲外は零）。
しかし各整曲線は $X$ の整成分であり、定理\ref{thm:point-graph-curve-degree}により
両次数は等しい。重複度を掛けて加えても等しいので $a=b$ を要する。

一方の文字の $d$ 次元真部分空間に支持される関係数は、
$(a,b)=(4-d,4)$ を用いると高々 $d$ となる。
特に固定因子を共有する二つの純関係はない。
相異なる純関係 $u_i\otimes v_i$ のうち四個以下で $u_i$ が従属とし、極小従属部分を取る。
その包の次元は個数より一つ小さいから、その純関係自身にも従属が必要である。
しかし $u_i$ の唯一の従属は全係数が非零であり、テンソルの従属を各 $v_i$ 側の汎関数で
評価すると、全 $v_i$ が比例することになる。固定因子の上界に反する。右側も同じである。
五つの相異なる純関係があれば、辺の各点対はそれぞれの第一因子または第二因子を消す。
四つを同時に消せないので、五添字を二個と三個に分け、一方を左、残りを右で消す有限個の
割当ての和集合に点対応が入る。各割当ての場所は直線と一点の積またはその逆である。
ここにグラフの整曲線は入れず、一次元性に反する。
\end{proof}

\begin{proposition}\label{prop:curve-graph-four-pure-frame}
四つの純関係がある場合、両文字の基底を取ればそれらは $u_i\otimes v_i$ ($1\le i\le4$) である。
二つの残余関係を対角項なしの二行列で表示する。
一次元の純曲線部分は被約な座標直線だけからなり、
各座標対 $I$ の左の直線を補対 $I^c$ の右の直線に射影同型で対応させる。
その直線が現れる必要十分条件は、両行列の $I\times I^c$ の二行二列ブロックが
一つの非退化行列を張ることである。
現れる左座標対のグラフは、一本の辺、交わる二辺、離れた二辺、三辺の道、または四辺の閉路のいずれかである。
対応するグラフの有限長部分（孤立部分と埋込み部分）の長さは順に $14,10,8,6,4$ となる。
\end{proposition}
\begin{proof}
四つの対角関係が消える場所は、座標集合を左右に分けた直積たちの和である。
各整曲線は両側に非定数なので、一点を含む分割は使えず、二対二の分割に入る。
その正方形は $\mathbb P^1\times\mathbb P^1$ で、残る二式はともに双次数 $(1,1)$。
グラフ成分を切る条件はちょうど表示した階数条件である。二式とも零または
因数分解する曲線全体が含まれることは一点ファイバー条件に反する。
この曲線の一般点では二つずつの非零座標が単元なので、四つの対角式が残り四座標を
スキームとして零にする。従って重複度は1で、有限台の部分を除いた純曲線は被約である。

ある頂点に三辺が入れば、その左座標点の像は三つの相補直線の共通点でなければならず、そのような点はない。
右側でも同様である。$n$ 本の左辺での頂点次数を $d_i$ とすれば相補グラフの頂点次数は $n-d_i$。
全て $d_i,n-d_i\le2$ なので非空グラフは列挙した五型に限られる。
交わる座標直線は異なる接方向を持つ通常二重点である。よって純曲線の算術種数は順に
$0,0,-1,0,1$ であり、定理\ref{thm:pure-curve-defect}の $18-4n+2g$ を代入して長さを得る。
この列挙は正則窓における各型の実在を主張するものではない。
\end{proof}
\path{computations/verify_curve_graph_pure_frames.py}は両次数の等号と相補座標グラフの全列挙を検算する。
